\chapter{Intelligence at the Interface}

The previous chapter described resistance. This chapter follows the disturbance that meets it. When a model begins to send commands, allocate work, or initiate actions, its outputs become inputs to a system with dynamics of its own.

A sentence can be evaluated as a sentence. A command must also be evaluated through the movement, load, feedback, and consequences it produces. The interface changes the task.

\section{From Solver to Driver}

\subsection{A Different Position in the Loop}

\subhead{A solution that has to act}

A solver compares possibilities and returns a candidate. A driver participates in changing the state of the world. It sends a signal through an actuator or an organization, observes a result, and helps determine the next action.

The phrase “phase change” describes this shift in role. The new system is evaluated on whether the transition occurs, remains stable, and can be recovered from when it goes wrong.

Ashby's law of requisite variety gives an early account of why the shift matters.\cite{ashby1956cybernetics} A regulator needs enough available distinctions and responses to counter the disturbances relevant to its task. Many words about a disturbance do not necessarily supply an actuator that can correct it.

A robot may describe a dropped object without being able to reach it in time. A logistics system may identify a delay without having a spare vehicle. Control authority has to exist at the location and time where the correction is needed.

\subhead{The answer becomes part of a cycle}

In a closed loop, the sequence is observation, decision, action, response, and a new observation. Each has latency and error. The model acts on an estimate of a state that may already have changed.

Wiener's work grew from precisely such prediction-and-action problems.\cite{wiener1948} Modern control theory evaluates a controller together with the plant, sensing, delay, and disturbance.\cite{astrom2008} The quality of a stand-alone forecast is one part of the result.

Supply-chain experiments show how people can also destabilize delayed feedback through their ordering decisions.\cite{sterman1989,lee1997} A system that reacts vigorously to a lagging signal can amplify the condition it is trying to correct.

Faster observation and computation can improve control. Their effect depends on the loop into which they enter. The gain has to be placed where the receiving process can use it.

\subsection{Four Conditions of Executability}

\subhead{Power, bandwidth, error, and safety}

An executable plan must fit four kinds of constraint.

First, the action needs adequate power and force at the relevant scale. Second, the receiving medium must be able to respond over the required bandwidth. Third, estimation, prediction, and execution errors must remain within the task's tolerances. Fourth, the failure modes must fit a defensible safety arrangement.

These conditions interact. A lighter actuator may move faster and tolerate less load. A higher gain may improve nominal tracking and excite an unmodeled mode. More sensors can improve observability while adding communications and maintenance requirements.

Safety engineering examines the interaction among components, people, organizations, and operating conditions.\cite{perrow1984,leveson2011} A component's success rate cannot establish the safety of every system into which it can be connected.

\subhead{Why the chosen path bends}

A shortest route on a map can have poor traction, a closed gate, an unsafe crossing, or an inaccessible entrance. A nominally optimal manufacturing schedule can eliminate the time needed for cleaning or repair. The path bends around conditions omitted from the first optimization.

The useful objective becomes a route that can be sustained at an acceptable total cost. It can include redundancy, staged release, human review, or spare capacity. These are design variables with measurable functions.

The receiving organization matters as much as the actuator. A clinical recommendation must enter a workflow with the information, authority, staffing, and follow-up necessary to act on it. An instruction no one can safely carry out remains an instruction.

\subhead{Embodiment as integration}

Brooks and later work on embodiment emphasized that useful behavior depends on the coupled body and environment.\cite{brooks1991,pfeifer2006} Sensors, morphology, compliance, feedback, and physical contact can perform work that a detached planner would otherwise have to specify.

Vision-language-action models extend this integration. RT-2 connected web-scale visual and language representations with robot action data.\cite{zitkovich2023} Its significance is the attempt to carry knowledge into action, where trajectories and outcomes impose another test.

Contemporary deployment shows the surrounding work. GXO's October 26, 2026 announcement expanded its relationship with Agility Robotics from an earlier live warehouse deployment and described planned additional sites.\cite{gxo2026digit} The relevant unit of progress is a robot within a workflow, with material handoffs, uptime, service, and supervision attached.

Autonomous-driving validation studies similarly examine the distance between observed operation and a sufficiently strong safety claim.\cite{kalra2016,koopman2016} Rare events, scenario coverage, operational limits, and the evidence supporting a release all matter.

A usable driver is an arrangement of model, body or tools, environment, and organization. Its capability is distributed across that arrangement.
